\chapter{Complexity — Classes}
%%% generated by the blueprint pipeline from TCSlib.Complexity.ClassOracle.Classes
\section{Overview}

Declarations of the Lean module \texttt{TCSlib.Complexity.ClassOracle.Classes}.

\section{Declarations}

\begin{definition}[Deterministic oracle time class]
\lean{Complexity.DTIMEOracle}
\label{Complexity.DTIMEOracle}
\leanok
\uses{Turing.FinOracleTM, Turing.FinOracleTM.DecidesInTime}
For an oracle language $O$ and a time bound $T\colon\mathbb{N}\to\mathbb{N}$, the
deterministic oracle time class consists of all languages $L$ for which there are a
natural number $c$ and a deterministic oracle Turing machine that decides $L$ relative
to $O$ within time $cT(n)$ on inputs of length $n$.
\end{definition}

\begin{definition}[Nondeterministic oracle time class]
\lean{Complexity.NTIMEOracle}
\label{Complexity.NTIMEOracle}
\leanok
\uses{Turing.FinOracleNDTM, Turing.FinOracleNDTM.DecidesInTime}
For an oracle language $O$ and a time bound $T\colon\mathbb{N}\to\mathbb{N}$, the
nondeterministic oracle time class consists of all languages $L$ for which there are a
natural number $c$ and a nondeterministic oracle Turing machine that decides $L$
relative to $O$ within time $cT(n)$ on inputs of length $n$.
\end{definition}

\begin{definition}[Polynomial-time oracle class]
\lean{Complexity.POracle}
\statementsource{BGS75}{
  pdf-2459e3840724-p002-b003,
  pdf-2459e3840724-p002-b004
}
\label{Complexity.POracle}
\leanok
\uses{Complexity.DTIMEOracle}
For an oracle language $O$, the polynomial-time oracle class is the union, over all
natural numbers $c$, of the deterministic oracle time classes with time bound $n^c+1$.
\end{definition}

\begin{definition}[Nondeterministic polynomial-time oracle class]
\lean{Complexity.NPOracle}
\statementsource{BGS75}{
  pdf-2459e3840724-p002-b003,
  pdf-2459e3840724-p002-b004
}
\label{Complexity.NPOracle}
\leanok
\uses{Complexity.NTIMEOracle}
For an oracle language $O$, the nondeterministic polynomial-time oracle class is the
union, over all natural numbers $c$, of the nondeterministic oracle time classes with
time bound $n^c+1$.
\end{definition}

\begin{theorem}[Polynomial time with any oracle]
\lean{Complexity.P_subset_POracle}
\label{Complexity.P_subset_POracle}
\leanok
\uses{Complexity.P, Complexity.POracle, Turing.FinTM.toFinOracleTM, Turing.FinTM.toFinOracleTM_computesInTime}
\difficulty{3}
For every binary language $O$ used as an oracle, every binary language decidable in
polynomial time belongs to the polynomial-time class with oracle $O$: $\mathsf{P}
\subseteq \mathsf{P}^{O}$.
\end{theorem}

\begin{theorem}[Deterministic oracle time is contained in nondeterministic oracle time]
\lean{Complexity.POracle_subset_NPOracle}
\label{Complexity.POracle_subset_NPOracle}
\leanok
\uses{Complexity.NPOracle, Complexity.POracle, Turing.FinOracleTM.toFinOracleNDTM, Turing.MultiTapeTM.indicator, Turing.OracleNDTM.initCfg, Turing.OracleNDTM.runWith, Turing.OracleTM.initCfg, Turing.OracleTM.runFrom, Turing.OracleTM.toOracleNDTM_runWith}
\difficulty{5}
For every language $O$ over $\{0,1\}$, every language decidable in polynomial time by a
deterministic Turing machine with oracle access to $O$ is also decidable in polynomial
time by a nondeterministic Turing machine with the same oracle; equivalently,
$	ext{P}^O\subseteq\text{NP}^O$.
\end{theorem}

\begin{definition}[Oracle machine constructed from a finite machine]
\lean{Complexity.queryTM}
\statementsource{BGS75}{
  pdf-2459e3840724-p002-b001,
  pdf-2459e3840724-p002-b002
}
\label{Complexity.queryTM}
\leanok
\uses{Turing.FinOracleTM, Turing.FinTM, Turing.OracleTM.answerAction, Turing.captureAction}
Given a finite Turing machine over the Boolean alphabet, with $k$ work tapes and state
set $Q$, form an oracle Turing machine with $k$ ordinary work tapes and one additional
query tape. Its state set is the disjoint union of $Q$ and three new states: a query
state, a positive-answer state, and a negative-answer state. The initial state is the
original machine’s initial state; transitions in states from $Q$ are obtained from the
original machine’s transitions by the specified action-capturing construction, while the
positive- and negative-answer states enact the corresponding oracle-answer steps and the
query state has no successor transition. The three distinguished states are pairwise
distinct.
\end{definition}

\begin{theorem}[Transition in the Positive-Answer State]
\lean{Complexity.queryTM_yes}
\label{Complexity.queryTM_yes}
\leanok
\uses{Complexity.queryTM, Turing.FinTM, Turing.OracleTM.answerAction}
\difficulty{1}
Let $F$ be a finite-state binary Turing machine, let $\mathrm{inp}$ be an optional
Boolean input symbol, and let $w$ assign a possibly blank Boolean symbol to each of the
constructed oracle machine's $k+1$ work tapes. In its positive-answer state, the
constructed machine's transition is the action that emits true, leaves the input and
work tapes unmoved and unwritten, and halts.
\end{theorem}

\begin{theorem}[Transition in the Negative-Answer State]
\lean{Complexity.queryTM_no}
\label{Complexity.queryTM_no}
\leanok
\uses{Complexity.queryTM, Turing.FinTM, Turing.OracleTM.answerAction}
\difficulty{1}
Let $F$ be a finite-state binary Turing machine, let $\mathrm{inp}$ be an optional
Boolean input symbol, and let $w$ assign a possibly blank Boolean symbol to each of the
constructed oracle machine's $k+1$ work tapes. In its negative-answer state, the
constructed machine's transition is the action that emits false, leaves the input and
work tapes unmoved and unwritten, and halts.
\end{theorem}

\begin{theorem}[Transition on an Embedded Original State]
\lean{Complexity.queryTM_agree}
\label{Complexity.queryTM_agree}
\leanok
\uses{Complexity.queryTM, Turing.FinTM, Turing.captureAction}
\difficulty{1}
Let $F$ be a finite-state binary Turing machine, let $s$ be one of its states, let
$\mathrm{inp}\in\{\mathrm{none},\mathrm{some}(b):b\in\{0,1\}\}$, and let $w$ assign a
possibly blank symbol to each of the $k+1$ work tapes of the constructed oracle machine.
At the state corresponding to $s$, its transition on $\mathrm{inp}$ and $w$ is the
capture transform of $F$'s transition on $\mathrm{inp}$ and the symbols on the first $k$
tapes, using the oracle machine's query state as the return state.
\end{theorem}

\begin{theorem}[Initial Configuration of the Constructed Oracle Machine]
\lean{Complexity.queryTM_init}
\label{Complexity.queryTM_init}
\leanok
\uses{Complexity.queryTM, Turing.FinTM, Turing.MultiTapeTM.initCfg, Turing.OracleTM.initCfg, Turing.captureCfg}
\difficulty{4}
Let $F$ be a finite-state binary Turing machine and let $x$ be a Boolean word. The
initial configuration of the oracle machine constructed from $F$ on input $x$ is the
capture image of $F$'s initial configuration, using the inclusion of $F$'s states into
the enlarged state set, the query state as the return state, and empty prefix and output
words.
\end{theorem}

\begin{theorem}[Polynomial reductions to an oracle class]
\lean{Complexity.mem_POracle_of_polyTimeReducible}
\label{Complexity.mem_POracle_of_polyTimeReducible}
\leanok
\uses{Complexity.POracle, Complexity.queryTM, Complexity.queryTM_agree, Complexity.queryTM_init, Complexity.queryTM_no, Complexity.queryTM_yes, Complexity.succ_pow_le, Turing.FinOracleTM.ComputesInTime, Turing.FinOracleTM.ComputesInTime.mono, Turing.FinTM.computesInTime_iff, Turing.MultiTapeTM.indicator, Turing.MultiTapeTM.initCfg, Turing.MultiTapeTM.runFrom, Turing.OracleTM.capture_reaches_qQuery, Turing.OracleTM.initCfg, Turing.OracleTM.queryString, Turing.OracleTM.queryString_bufferTape, Turing.OracleTM.queryTapeIdx, Turing.OracleTM.runFrom, Turing.OracleTM.runFrom_add, Turing.OracleTM.runFrom_two_of_qQuery, Turing.captureCfg}
\difficulty{6}
Let $L$ and $O$ be languages over the binary alphabet. If $L$ is polynomial-time
many-one reducible to $O$, then $L$ is decidable in polynomial time by a Turing machine
with oracle access to $O$.
\end{theorem}

\begin{theorem}[An oracle language is decidable with its oracle]
\lean{Complexity.oracle_mem_POracle}
\label{Complexity.oracle_mem_POracle}
\leanok
\uses{Complexity.POracle, Complexity.PolyTimeReducible.refl, Complexity.mem_POracle_of_polyTimeReducible}
\difficulty{1}
For every language $O\subseteq\{0,1\}^*$, the language $O$ belongs to the
polynomial-time class relative to oracle $O$.
\end{theorem}

\begin{definition}[Negating machine output]
\lean{Complexity.negOut}
\label{Complexity.negOut}
\leanok
\uses{Turing.FinOracleTM}
Given a finite Boolean oracle Turing machine, form a finite Boolean oracle Turing
machine with the same number of ordinary work tapes, states, initial state, query state,
and answer states. For every transition, preserve its input-head movement, work-tape
actions, and successor state, and replace each emitted output bit by its Boolean
negation; the three distinguished states remain pairwise distinct.
\end{definition}

\begin{definition}[Negating configuration output]
\lean{Complexity.negCfg}
\label{Complexity.negCfg}
\leanok
\uses{Turing.Cfg}
Given a configuration of a machine with $k+1$ work tapes, Boolean work-tape symbols, a
state set $S$, and a fixed Boolean input string $x$, form a configuration with the same
state, input-head position, work-tape contents, and work-tape head positions, but
replace every bit in the output string by its Boolean negation.
\end{definition}

\begin{theorem}[One-step output complementation]
\lean{Complexity.negOut_step}
\label{Complexity.negOut_step}
\leanok
\uses{Complexity.negCfg, Complexity.negOut, Turing.Cfg, Turing.Cfg.inputSymbol, Turing.Cfg.workTapeSymbols, Turing.FinOracleTM, Turing.OracleTM.step}
\difficulty{5}
Let $M$ be a finite-state oracle Turing machine over Boolean symbols, let $O$ be an
oracle language, let $x$ be an input bit string, and let $c$ be a configuration of $M$
on $x$. One step of the machine whose ordinary transition outputs are complemented,
starting from the configuration whose existing output bits are complemented, equals the
configuration obtained by taking one step of $M$ with oracle $O$ from $c$ and then
complementing every output bit.
\end{theorem}

\begin{theorem}[Output complementation along a run]
\lean{Complexity.negOut_runFrom}
\label{Complexity.negOut_runFrom}
\leanok
\uses{Complexity.negCfg, Complexity.negOut, Complexity.negOut_step, Turing.Cfg, Turing.FinOracleTM, Turing.OracleTM.runFrom, Turing.OracleTM.runFrom_comm_of_step}
\difficulty{1}
Let $M$ be a finite-state oracle Turing machine over Boolean symbols, let $O$ be an
oracle language, let $x$ be an input bit string, and let $c$ be a configuration of $M$
on $x$. For every natural number $t$, running the machine with each transition’s output
bit complemented for $t$ steps from the configuration whose existing output bits have
all been complemented gives exactly the configuration obtained by first running $M$ with
oracle $O$ for $t$ steps and then complementing every output bit.
\end{theorem}

\begin{theorem}[Complement closure of polynomial oracle time]
\lean{Complexity.compl_mem_POracle}
\label{Complexity.compl_mem_POracle}
\leanok
\uses{Complexity.POracle, Complexity.negCfg, Complexity.negOut, Complexity.negOut_runFrom, Turing.MultiTapeTM.indicator, Turing.OracleTM.initCfg, Turing.OracleTM.runFrom}
\difficulty{4}
For any binary oracle language $O$ and any binary language $L$, if $L$ is decidable in
polynomial time by a finite oracle Turing machine with oracle $O$, then its complement
is also decidable in polynomial time with the same oracle.
\end{theorem}

\begin{theorem}[Polynomial growth estimate]
\lean{Complexity.b1_poly_succ_le}
\label{Complexity.b1_poly_succ_le}
\leanok
\difficulty{3}
For all natural numbers $c,k,n$,\[c(n^k+1)+1\le (2c+1)(n+1)^k.\]
\end{theorem}

\begin{theorem}[Polynomial budget bound]
\lean{Complexity.b1_budget}
\label{Complexity.b1_budget}
\leanok
\uses{Complexity.b1_poly_succ_le, Complexity.succ_pow_le}
\difficulty{4}
For all natural numbers
$a,c,k,C,d,n$,\[a\bigl(c(n^k+1)+1\bigr)\bigl(c(n^k+1)+1+C(c(n^k+1)+1)^d\bigr)\le
a(1+C)(2c+1)^{d+2}2^{k(d+2)}\bigl(n^{k(d+2)}+1\bigr).\]
\end{theorem}

\begin{theorem}[Polynomial oracle collapses to ordinary polynomial time]
\lean{Complexity.POracle_eq_P_of_mem_P}
\label{Complexity.POracle_eq_P_of_mem_P}
\leanok
\uses{Complexity.DTIME, Complexity.P, Complexity.POracle, Complexity.P_subset_POracle, Complexity.b1_budget, Complexity.mem_P_iff, Complexity.mem_P_of_dtime_le, Turing.FinOracleTM.exists_plain_decider, Turing.FinTM.ComputesInTime.mono}
\difficulty{4}
Let $O$ be a language over $\\{0,1\\}$. If $O$ is decidable in polynomial time, then the
class of languages decidable in polynomial time with oracle access to $O$ is exactly the
ordinary polynomial-time class: $	ext{P}^O=	ext{P}$.
\end{theorem}
